\documentclass[output=paper]{langscibook}
\ChapterDOI{10.5281/zenodo.22880302}
\title{\textit{Ana suva} – Fish poisoning roots}
\author{Helen Kobaa Magum, Shalom Magum and Joyce Maion}
\abstract{}
\begin{document}
\maketitle\label{chap:37}
\begin{xltabular}{\textwidth}{lXX}
	(\ref{ex:text30-1})&\textit{Enaa kahi vahutate ni nom havee to suvasuva vaavaha vo raara,}&I am going to describe how we poison fish,\\
	(\ref{ex:text30-2})&\textit{beara suvasuva nio ana suva.}&when we poison fish with the \textit{suva} root, the poisonous root.\\
	(\ref{ex:text30-3})&\textit{Beara tau nao tea suvasuva, eara toro matatopo amaa taba to pasi dee raara.}&Before we go to poison fish, we must prepare the things that we will take there.\\
	(\ref{ex:text30-4})&\textit{Momohu eara na tagihu rara rasuu,}&First we go into the forest.\\
	(\ref{ex:text30-5})&\textit{Eara repaa kasi amaa ana suva en to neba komana raara e.}&We dig out the \textit{suva} roots where we plant them.\\
	(\ref{ex:text30-6})&\textit{Eara repaa dee voosu vaan,}&Then we carry them home to the village,\\
	(\ref{ex:text30-7})&\textit{eara repaa ihuana koa rara vaan beara sisira rara o tarahi vaa tea roava ge tea bon.}&and just wait in the village, while we observe the ebbing tide during the day or the night.\\
	(\ref{ex:text30-8})&\textit{Beara tara voen, a masi na mataa nana tea suvasuva,}&When we see that the reef is good for fish poisoning,\\
	(\ref{ex:text30-9})&\textit{eara repaa vaobete a ana suva teo hoi, bara a vasu.}&we put the \textit{suva} roots into a basket, and a stone.\\
	(\ref{ex:text30-10})&\textit{Eara repaa kae eve vo tea masi.}&Then we carry it to the reef.\\
	(\ref{ex:text30-11})&\textit{Beara suguna tea kopua to pasi suva ni raara,}&When we arrive at the pool that we intend to poison,\\
	(\ref{ex:text30-12})&\textit{eara repaa vahio o hoi ana suva.}&we put the basket with the \textit{suva} roots down.\\
	(\ref{ex:text30-13})&\textit{Eara repaa taneo tea tasu varavuravu a ana suva bona vasu.}&Then we start beating and squashing the \textit{suva} roots with the stone.\\
	(\ref{ex:text30-14})&\textit{Tasu batara, eara repaa peo komana kopua, bara amaa vavaarihi nae.}&We beat them and then scatter them in the pool and along its sides.\\
	(\ref{ex:text30-15})&\textit{Saka o peho ana haana to tasu raara,}&It is not a single root that we beat,\\
	(\ref{ex:text30-16})&\textit{a bero, ore suguna komana beo iana bubuutoo tavusu maa tea maa avaava vasu vaa komana kopua.}&it is many, until the fish struggle out of the rocky caves in the pool.\\
	(\ref{ex:text30-17})&\textit{Beara tara voen, eara repaa antee tea tasu a ana suva.}&When we see this, we stop beating the \textit{suva} roots.\\
	(\ref{ex:text30-18})&\textit{Eara repaa avete o iana vai to paa mateu.}&Then we collect the fish that have died.\\
	(\ref{ex:text30-19})&\textit{Beara avete vakavara, eara repaa tara voen o iana paa kavarau,}&When we have finished collecting the fish and see that the fish are finished,\\
	(\ref{ex:text30-20})&\textit{eara repaa nao pete tea rake ta meha kopua.}&we move on to search for another pool.\\
\end{xltabular}
\end{document}